% paper2_discussion.tex -- Discussion and Conclusion of Paper 2, input by
% papers/paper2_nitrogen.tex.
\section{Discussion}

The response of maize to nitrogen followed the familiar pattern of
diminishing returns, and the optimum rate differed between sites in line with
their yield potential: it was highest at the temperate site, where water was
least limiting, and lowest at the semi-arid site, where yield was capped by
water rather than by nitrogen. A single recommended rate across the three sites
would therefore have over-fertilised the dry site and under-fertilised the
temperate one, which is why site-specific recommendations are central to
improving nitrogen use efficiency \parencite{cassman2002}.

Agronomic and recovery efficiency both fell as the rate increased. This is
expected, because each additional kilogram of nitrogen is applied to a crop
whose demand is already closer to being met. The practical consequence is a
trade-off: rates near the agronomic optimum maximise yield, but the last
increments of nitrogen are used least efficiently and contribute most to the
surplus at risk of loss. The recovery values found here lie within the range
reported for cereal systems, where much of the applied nitrogen is not taken
up by the crop \parencite{ladha2005}.

The optimum estimated from a quadratic-plateau function is sensitive to the
form of the function and to the spread of rates tested, and it ignores the
prices of grain and fertiliser. An economic optimum, at which the value of the
extra grain equals the cost of the extra nitrogen, is always lower than the
agronomic optimum. Combining moderate rates of mineral nitrogen with organic
inputs, as integrated soil fertility management\index{integrated soil fertility management} recommends, is one way to
retain most of the yield benefit while improving efficiency
\parencite{vanlauwe2010}.

\section{Conclusion}

Maize yield rose with nitrogen to a plateau at all three sites, with agronomic
optimum rates of \res{p2-nopt-S1}, \res{p2-nopt-S2} and
\res{p2-nopt-S3}~\si{\kilo\gram\per\hectare} at S1, S2 and S3. Nitrogen use
efficiency declined steadily with rate, so the highest rates gave little extra
grain for the nitrogen applied. Recommendations should be set site by site and
weighed against efficiency as well as yield. The data are synthetic and the
values illustrate the calculation only.
